\section{Scalar inputs and the entropy--energy identity}\label{sec:scalar}

This section provides the deterministic coefficients for the proposal.
The main output is \eqref{eq:3}: it balances the proposal's information
cost through the matrix channel against its energy gain, up to a small
terminal error.
We derive it after stating the equilibrium inputs and controlling a
continuation of the scalar diffusion beyond the support of the Parisi
measure.

\subsection{Exact equilibrium inputs and normalization}

Let $\mu_\beta$ be a minimizing Parisi probability measure for
zero-field Gaussian SK at the fixed finite $\beta>1$. We use the
following external results: Guerra's finite-volume pressure upper
bound~\cite[Theorem~3]{Guerra2003}; spatial regularity and stability
under step approximation for the Parisi equation
\cite[Propositions~1--2]{AC2015}; and full interval support with scalar
stationarity~\cite[Theorem~1.1, Lemma~2.2, Proposition~2.4, and
Lemma~2.5]{Lopatto2026}. In particular,
\[
 \operatorname{supp}(\mu_\beta)=[0,q_*],\qquad 0<q_*<1.
\]
The cited support theorem is used only for this zero-field Gaussian
pair-interaction model and for each fixed $\beta>1$.

Use scalar variance time $v=\beta^2s$, and put
$\alpha_v=\mu_\beta([0,v/\beta^2])$ for $0\le v\le\beta^2$.
Let $\Phi$ solve the backward Parisi equation
\[
 \Phi_v=-\tfrac12(\Phi_{yy}+\alpha_v\Phi_y^2),
 \qquad \Phi(\beta^2,y)=\log(2\cosh y).
\]
Derivatives of $\Phi$ are indicated by subscripts. Equivalently, if
$u$ is the scalar solution in~\cite{Lopatto2026}, then
$\Phi(v,y)=u(v/\beta^2,y)+\log2$. Combining the source's independent
ordered pair couplings gives the unordered pair coefficient
$\beta/\sqrt n$ used here, plus a centered diagonal term. Thus its
scalar normalization agrees with the completed Hamiltonian.

Set $v_*=\beta^2q_*$. In these units the pressure functional and the
finite-volume upper bound read
\[
 \mathcal P=\Phi(0,0)-\frac1{2\beta^2}
                 \int_0^{\beta^2}v\alpha_v\,dv,
 \qquad \mathbb E F(G)\le n\mathcal P.
\]
The centered diagonal does not affect this expected pressure. Since
$\alpha_v=1$ above $v_*$, the Gaussian recursion gives
\[
 \Phi(v_*,y)=\log(2\cosh y)+\frac{\beta^2-v_*}{2}.
\]
Only the pressure upper bound is needed below.

Continue the endpoint formula for the potential to all $v>v_*$ by
\[
 \Phi(v,y)=\log(2\cosh y)+\frac{\beta^2-v}{2},\qquad \alpha_v=1.
\]
This agrees with the original solution up to $\beta^2$ and satisfies the
same PDE for arbitrarily large $v$. For every $v\ge0$, put
$m(v,y)=\Phi_y(v,y)$ and $a(v,y)=\Phi_{yy}(v,y)$; thus
$m(v,y)=\tanh y$ and $a(v,y)=1-\tanh^2y$ when $v\ge v_*$.
Let $Z$ be standard Brownian motion and solve
\[
 dY_v=\alpha_vm(v,Y_v)\,dv+dZ_v,\qquad Y_0=0.
\]
The drift is bounded and spatially Lipschitz. We abbreviate its evaluated
scalar functions by $m_v=m(v,Y_v)$ and $a_v=a(v,Y_v)$. The cited
stationarity identity, the differentiated PDE, and It\^o's formula give
\[
 m_0=0,\qquad dm_v=a_v\,dZ_v,\qquad
 \mathbb E m_v^2=v/\beta^2\quad(0\le v\le v_*).
\]
Full interval support is what makes the last identity available for
every time in $[0,v_*]$; the variational identity on the support is
also given in~\cite[Proposition~1.1 and Corollary~3.10]{JT2017}.
The martingale equation continues beyond $v_*$ by direct It\^o
calculus. Spatial derivatives are continuous under step approximation,
so none of these assertions requires time differentiability at a CDF
atom.

We will repeatedly use
\[
 0<a(v,y)\le1-m(v,y)^2,\qquad |m(v,y)|<1.
\]
Here is the derivative bound in a form suited to step approximation.
One step with CDF height $\lambda\in[0,1]$ and variance $\delta>0$
takes a profile $f$ to
\[
 (Tf)(y)=\lambda^{-1}\log\mathbb E
             e^{\lambda f(y+\sqrt\delta U)},\qquad U\sim N(0,1),
\]
with $Tf(y)=\mathbb E f(y+\sqrt\delta U)$ when $\lambda=0$.
Write $\mathbb E_\lambda$ for the Gaussian expectation tilted by
$e^{\lambda f(y+\sqrt\delta U)}$ and normalized to have total mass one.
Differentiating gives
\[
 (Tf)'=\mathbb E_\lambda f',\qquad
 (Tf)''=\mathbb E_\lambda f''+
                 \lambda\operatorname{Var}_\lambda(f').
\]
If $0<f''\le1-(f')^2$, these formulas imply
\[
 0<(Tf)''\le1-\bigl((Tf)'\bigr)^2
              -(1-\lambda)\operatorname{Var}_\lambda(f')
 \le1-\bigl((Tf)'\bigr)^2.
\]
Inducting from $f(y)=\log(2\cosh y)$ proves the bound for step CDFs;
scalar stability passes the non-strict bound to the minimizing CDF.
The positive lower bound for $\Phi_{yy}$ in
\cite[Proposition~2.4]{Lopatto2026} supplies strict positivity in the
limit. Hence $m$ is strictly increasing and bounded by one in absolute
value, so $|m|<1$. On the extension these assertions follow from $\tanh$.

\subsection{Tail estimates for the extended diffusion}

Define the deterministic functions
\[
 p(v)=\mathbb E a_v,\qquad r(v)=\mathbb E a_v^2,
 \qquad \rho(v)=\beta\sqrt{r(v)}.
\]
The martingale second-moment identity and continuity give
$r(v)=1/\beta^2$ on $[0,v_*]$, including the endpoints.
The next lemma controls the continuation and shows that $\rho$ is
an admissible noise-mixing coefficient.

\begin{lemma}[Scalar tail bounds]\label{lem:scalar-tail}
For $v\ge v_*$, the extended diffusion satisfies
\[
 \begin{gathered}
 \rho(v)\le1,\qquad -2\le(\log r)'(v)\le0,\\
 p(v)=1-\mathbb E m_v^2,\qquad p'(v)=-r(v),
 \qquad p(v)\le C_\beta e^{-v/2}.
 \end{gathered}
 \tag{1}\label{eq:1}
\]
At the endpoint derivatives are understood from the right.
\end{lemma}

\begin{proof}
We use the transformed-density method of
\cite[Proposition~7.1]{Lopatto2026} and the moment comparison in
\cite[Proposition~9.3]{Lopatto2026}, giving the required estimates for
the arbitrary-time continuation.
We first prove two moment ratios. For $v>v_*$, the density of $Y_v$
divided by $\cosh y$ is symmetric and nonincreasing in $|y|$; it is
enough to establish the resulting moment inequalities for step CDFs
and then pass to the limit. To verify the shape property, take a step
CDF with mass $\alpha^{(j)}$ on $[v_j,v_{j+1}]$, and write
$\varphi_t$ for the centered Gaussian density of variance $t$.
The diffusion transition density on that interval is
\[
 K_j(z,y)=\varphi_{v_{j+1}-v_j}(y-z)
 \exp\{\alpha^{(j)}[\Phi(v_{j+1},y)-\Phi(v_j,z)]\}.
\]
The Gaussian recursion normalizes this kernel; its Doob transform has
drift $\alpha^{(j)}\Phi_y$. In a product of these kernels, the
intermediate factors at $y_j$ reduce to
\[
 \exp\{-(\alpha^{(j)}-\alpha^{(j-1)})\Phi(v_j,y_j)\}.
\]
They are even and nonincreasing in $|y_j|$, because the masses are
nondecreasing and $\Phi$ is even and convex. The first positive-time
Gaussian convolution turns the initial point mass at zero into a
symmetric decreasing density. Multiplication by the displayed factors
and further Gaussian convolutions preserve this property. For
convolution, this follows by representing the profiles through their
centered-interval level sets. After extending the final mass to one,
the endpoint factor is a constant times $\cosh y$. Dividing by this
factor leaves the asserted symmetric decreasing profile.

Weighting either the density proportional to $\operatorname{sech}y$
or that proportional to $\operatorname{sech}^3y$ by this decreasing
profile can only increase the expectation of
$\operatorname{sech}^2y$: two nonincreasing functions of $|y|$ have
nonnegative covariance. Integration by parts gives
\[
 \frac{\int_{\mathbb R}\operatorname{sech}^{k+2}y\,dy}
      {\int_{\mathbb R}\operatorname{sech}^{k}y\,dy}
 =\frac{k}{k+1}\qquad(k>0).
\]
Consequently $r/p\ge1/2$ and $\mathbb E a_v^3/r(v)\ge3/4$.
The passage from step CDFs is justified by scalar stability.

On the extension $a_v=1-m_v^2$, and direct It\^o calculations yield
\[
 p'=-r,\qquad r'=4r-6\mathbb E a_v^3.
\]
The two moment ratios imply $p'\le-p/2$ and $r'\le0$.
Since $0<a_v\le1$, the second identity also gives $r'\ge-2r$.
Integrating the first differential inequality proves the exponential
bound on $p$. Finally $r(v_*)=1/\beta^2$ and monotonicity imply
$\rho\le1$.
\end{proof}

\subsection{Balancing scalar entropy and energy}

Let $\mathcal H(z)$ be the entropy of a sign with mean $z$:
\[
 \mathcal H(z)=-\frac{1+z}{2}\log\frac{1+z}{2}
              -\frac{1-z}{2}\log\frac{1-z}{2},
\]
with the usual continuous values at $z=\pm1$. Define, for $V>v_*$,
\[
 e(V)=\int_0^V\sqrt{r(v)}\,p(v)\rho(v)\,dv,
 \qquad j(V)=\frac12\int_0^V\rho(v)^2
                    \mathbb E\frac{a_v^2}{1-m_v^2}\,dv.
\]
The quantity $e(V)$ will be the energy per spin of the matrix proposal;
$j(V)$ will measure the information acquired through its matrix channel.
The independent noise introduced later accounts for the factor
$\rho(v)^2$ in this information cost.

Up to $v_*$, It\^o's formula, the PDE and stationarity give
\[
 \begin{aligned}
 \mathbb E[Y_{v_*}m_{v_*}]
   &=\int_0^{v_*}\bigl(p(v)+\alpha_v v/\beta^2\bigr)\,dv,\\
 \mathbb E\Phi(v_*,Y_{v_*})
   &=\Phi(0,0)+\frac12\int_0^{v_*}\alpha_v v/\beta^2\,dv.
 \end{aligned}
\]
These formulas also follow by step approximation with bounded first
and second spatial derivatives. Substitute the endpoint value of
$\Phi$ and use
$\mathcal H(\tanh y)=\log(2\cosh y)-y\tanh y$ to obtain
\[
 \int_0^{v_*}p(v)\,dv+\mathbb E\mathcal H(m_{v_*})
 =\mathcal P-\frac{\beta^2}{4}(1-q_*)^2.
 \tag{2}\label{eq:2}
\]

\begin{lemma}[Entropy--energy identity]\label{lem:scalar-balance}
For every $V>v_*$,
\[
 \log2-j(V)+\beta e(V)
 =\mathcal P-\frac{\beta^2}{4}p(V)^2.
 \tag{3}\label{eq:3}
\]
\end{lemma}

\begin{proof}
The same martingale calculation applied to binary entropy yields
\[
 \frac12\int_0^{v_*}\mathbb E\frac{a_v^2}{1-m_v^2}\,dv
 =\log2-\mathbb E\mathcal H(m_{v_*}).
\]
Localization justifies It\^o's formula near $\pm1$; the integrand is
bounded because $a\le1-m^2$. As $\rho=1$ and
$\beta\sqrt r=1$ on $[0,v_*]$, these identities prove
\eqref{eq:3} at $V=v_*$. Above $v_*$, the derivative of its left
side is $(\beta^2/2)r(v)p(v)$. Indeed $a_v=1-m_v^2$ there.
The derivative of $-\beta^2p(v)^2/4$ is the same, by $p'=-r$.
Integration proves the claim for every $V>v_*$.
\end{proof}

The terminal error in \eqref{eq:3} is at most $C_\beta e^{-V}$.
We will take $V$ proportional to $\log n$, so that this error is a
negative power of the dimension while the scalar spin is nearly pinned.
